\begin{helper}
Let \(\eta,\varepsilon_1\in\mathbb R\) and \(n\in\mathbb N\).  Put
\[
L=\log n,\quad
k_R=(1+\eta)\sqrt{nL},\quad
K=\noderef{TwoBiteNaturalIndependenceScale}(1+\eta,n),\quad k=(K:\mathbb R),
\]
\[
M=\noderef{TwoBiteNaturalReducedVertexCount}(n),\quad m=(M:\mathbb R),
\quad
p=\frac12\sqrt{\frac Ln},\quad
t_1=\frac{\sqrt{nL}}{\log L}.
\]
Assume \(0<\eta\), \(0<\varepsilon_1\), \(0<n\), \(0<L\), \(0<\log L\),
\(0\le k\), \(k_R\le k\),
\[
0\le 2pm,\qquad
2pm\le\frac{\sqrt n}{L^{3/2}},
\]
and
\[
\frac{2\log L}{(1+\eta)L^4}\le\varepsilon_1.
\]
Then
\[
\frac{2k}{t_1}\noderef{RealChooseTwo}(2pm)\le\varepsilon_1k^2.
\]
\end{helper}
\begin{proof}
Let \(x=2pm\).  Since \(x\ge0\), the upper bound in
\(\noderef{RealChooseTwoQuadraticBounds}\) gives
\[
\noderef{RealChooseTwo}(x)\le x^2.
\]
The mass estimate gives
\[
x^2\le \left(\frac{\sqrt n}{L^{3/2}}\right)^2.
\]
The factor \(2k/t_1\) is nonnegative, so
\[
\frac{2k}{t_1}\noderef{RealChooseTwo}(x)
\le
\frac{2k}{t_1}\left(\frac{\sqrt n}{L^{3/2}}\right)^2.
\]
Using \(t_1=\sqrt{nL}/\log L\), \(L^{3/2}=L\sqrt L\), and
\(k_R=(1+\eta)\sqrt{nL}\), the right side is
\[
k\frac{2\log L}{(1+\eta)L^4}k_R.
\]
The displayed threshold bounds this by \(k\varepsilon_1 k_R\).  Since
\(k_R\le k\), \(k\ge0\), and \(\varepsilon_1>0\),
\[
k\varepsilon_1 k_R\le k\varepsilon_1 k=\varepsilon_1 k^2.
\]
Combining the inequalities proves the claim.
\end{proof}
